\section{The modular kernel of an alternating Heisenberg carrier}
\label{sec:heisenberg-modular-kernel}

\paragraph{Status and provenance.}
D-TD20--21, reserved question Q2ei, is separate from D-TD17--19.
Root suggested restricting the symplectic label space to the kernel
of its alternating quotient and grouping dual isotypic components.
The proof below was developed by \texttt{/root/imprimitive}; it works
over the original field \(\mathbf F_p\), so no unproved descent
from an algebraic closure is needed. Root supplied the shorter
final total-component count. The shared derivation is recorded
explicitly. The complete author mathematical and dependency-interface
reading and complete root and independent whole-proof reviews pass. No finite
mathematical computation is needed or claimed.
No multiplicity-rank estimate or bound on the Heisenberg exponent is
claimed. The accepted orbit theorem D-TD16 is the source input.

\subsection*{D-TD20. Orthogonal blocks over the original finite field}

Let \(Q\) be a nonzero symplectic vector space of dimension \(2m\)
over \(\mathbf F_p\), including \(p=2\), and let
\[
 H\le\operatorname{Sp}(Q),\qquad
 Q\text{ irreducible as an }\mathbf F_pH\text{-module},
 \qquad N\triangleleft H.
 \tag{TDI.1}
\]
The restriction to \(N\) is not assumed to have order prime to \(p\).
A module is called homogeneous here if it is a direct sum of copies
of one irreducible \(\mathbf F_pN\)-module. It need not be absolutely
homogeneous over an algebraic closure.

\begin{lemma}[D-TD20]
The restriction \(Q|_N\) is semisimple. Let
\[
                     Q=Q_1\oplus\cdots\oplus Q_s
 \tag{TDI.2}
\]
be its canonical isotypic decomposition, indexed by the distinct
irreducible \(\mathbf F_pN\)-types occurring. The group \(H\)
permutes these components transitively, and its action on their set
kills \(N\). Duality of their irreducible types gives an involution
\(i\mapsto i^*\) with the following properties:
\begin{enumerate}
\item A component with \(i=i^*\) is nondegenerate for the restricted
alternating form.
\item If \(i\ne i^*\), the two components are individually totally
isotropic and pair perfectly with one another. Their sum is
nondegenerate.
\end{enumerate}
Consequently the orbits of this involution give a canonical orthogonal
sum of nonzero symplectic spaces over \(\mathbf F_p\),
\[
                    Q=B_1\perp\cdots\perp B_b,
                     \qquad b\le s\le2b,
 \tag{TDI.3}
\]
which \(H\) permutes transitively. All \(B_j\) have the same even
dimension.

Suppose additionally that \(Q\) is represented projectively
irreducibly on a complex space \(W\), with a nondegenerate scalar
commutator pairing, and that an actual finite projective normalizer image
\(P\subset\operatorname{PGL}(W)\) induces precisely \(H\) on \(Q\).
Then (TDI.3) induces an actual transitive orbit of full complex
matrix algebras and a tensor factorization
\[
       W\simeq W_1\otimes\cdots\otimes W_b,
                 \qquad \dim W_j=p^{\dim B_j/2}\ge2.
 \tag{TDI.4}
\]
Every operator in \(P\) permutes these embedded factor algebras.
The assertion also holds inside \(V=W\otimes R\) for any entire
external partner \(R\) and any actual group mapping onto \(P\).
It uses no splitting of a projective normalizer extension.
\end{lemma}

\begin{proof}
A nonzero finite-dimensional module has a simple submodule, by
choosing a nonzero submodule of minimal dimension. The sum of all
simple \(N\)-submodules of \(Q\), its \(N\)-socle, is nonzero.
Normality of \(N\) makes this sum \(H\)-invariant. Irreducibility
in (TDI.1) makes it equal to \(Q\).
This sum is semisimple: take a maximal direct sum of simple
submodules; a simple submodule not contained in it intersects
it trivially and would enlarge it. Finite dimension makes the
maximal sum contain every simple submodule. Thus (TDI.2) is the
usual canonical decomposition by simple isomorphism types, over
\(\mathbf F_p\) itself. No characteristic-zero Maschke assertion
has been applied in characteristic \(p\).

Conjugation by \(h\in H\) sends a simple \(N\)-submodule to one
with the conjugated \(N\)-type. It therefore permutes the whole
isotypic components. The sum over any orbit would be a nonzero
\(H\)-submodule, so there is just one orbit. Each component is
an \(N\)-submodule and hence is individually stabilized by \(N\).

The nondegenerate alternating form identifies \(Q\) with its
contragredient \(N\)-module. The dual of a simple finite-dimensional
module is simple: an invariant proper subspace of its dual would
have a nontrivial proper annihilator in the original module.
The form between \(Q_i\) and \(Q_k\) induces an \(N\)-map
from \(Q_i\) to \(Q_k^*\). Maps between distinct simple types
vanish, since a nonzero map between simples is an isomorphism.
It follows that
\[
                 \langle Q_i,Q_k\rangle=0
                         \quad\text{unless }k=i^*.
 \tag{TDI.5}
\]
Every dual type must occur, or the corresponding component would
be orthogonal to all of \(Q\). Thus duality is an involution on
the actual list of types. If it fixes \(i\), the radical of the
form on \(Q_i\) is orthogonal to every component by (TDI.5),
and is zero. If it exchanges \(i,i^*\), the form vanishes on
each component separately. A vector in either cross-pairing kernel
would again be orthogonal to all of \(Q\); both kernels are zero.
The cross-pairing is therefore perfect and their sum is nondegenerate.
This proves the orthogonal decomposition and the count in (TDI.3).
The form is alternating also at \(p=2\); every nonzero nondegenerate
block consequently has even positive dimension.

The action of \(H\) commutes with this dual-pair grouping because
it preserves the form. Its transitivity on the components gives
transitivity on the grouped blocks. It also gives equality of their
dimensions. This proves all finite-field assertions without any
extension of scalars or absolute-irreducibility assumption.

We verify that these are actual complex tensor factors. The scalar
commutator form is also an \(H\)-invariant nondegenerate alternating
form on \(Q\). Apply the same argument (TDI.5) to this form: its
orthogonal blocks are the same canonical groups of dual simple types,
and each is nondegenerate. Thus (TDI.3) holds for this form as well;
no assertion that two invariant forms are scalar multiples is needed.
For the rest of this proof use this commutator form. Choose unitary
Weyl representatives \(L_v\) of the projective labels \(v\in Q\).
The accepted intrinsic proof D-TD05 gives
\[
 \{L_v:v\in Q\}\text{ a basis of }\operatorname{End}_{\mathbf C}(W),
 \qquad
 L_vL_wL_v^{-1}L_w^{-1}
                   =\zeta^{\langle v,w\rangle}I,
 \tag{TDI.6}
\]
where \(\zeta\) is a primitive \(p\)-th root. Changing \(\zeta\)
only rescales this form and leaves orthogonality unchanged.
For each block put
\[
             \mathcal A_j=\operatorname{span}_{\mathbf C}
                                  \{L_v:v\in B_j\}.
 \tag{TDI.7}
\]
It is a unital algebra, since multiplying labels adds them.
Different \(\mathcal A_j\) commute by the orthogonality in (TDI.3).
Their dimensions are \(p^{\dim B_j}\) by the basis in (TDI.6).

For completeness, each \(\mathcal A_j\) is a full matrix algebra.
Take a symplectic basis of \(B_j\) and lifts \(X_1,\ldots,X_d\),
\(Z_1,\ldots,Z_d\), where \(2d=\dim B_j\). Scale each lift to
have order \(p\). The \(Z_i\) commute and are simultaneously
diagonalizable over \(\mathbf C\). Their common characters are
indexed by \(\mathbf F_p^d\). The commutator relations make the
\(X_i\) translate these characters independently and transitively;
starting from any nonzero common eigenspace shows that all \(p^d\)
characters occur with the same multiplicity. Choose a basis in one
such eigenspace and transport it by the commuting \(X_i\).
In these coordinates the generated algebra is
\(\operatorname{Mat}_{p^d}(\mathbf C)\otimes I\): the Fourier
projectors for the characters and the translations give its matrix
units. This argument remains valid for \(p=2\); only the alternating
commutator pairing is used, not a quadratic refinement.

The product of the commuting algebras \(\mathcal A_j\) contains
all \(L_v\), because \(v\) is the sum of its block components
and the corresponding product of lifts differs from \(L_v\)
only by a scalar. Hence multiplication gives a surjective map
\[
          \bigotimes_{j=1}^b\mathcal A_j
                       \longrightarrow\operatorname{End}(W).
 \tag{TDI.8}
\]
Both sides have dimension \(p^{2m}\), so this is an isomorphism.
The tensor product of the standard modules for these full matrix
algebras gives (TDI.4). Thus no multiplicity factor within \(W\)
has been discarded.

For an operator of \(P\) inducing \(h\in H\), conjugation sends
\(L_v\) to a nonzero scalar times \(L_{hv}\). It therefore sends
\(\mathcal A_j\) to the algebra belonging to \(hB_j\).
This proves actual normalization and permutation of the factor
algebras, independently of scalar choices. A group acting on
\(W\otimes R\) through this projective factor image has the same
conjugation action on \(\mathcal A_j\otimes I_R\); the entire
external commutant remains present. This proves the last assertion.
\end{proof}

\subsection*{D-TD21. Homogeneity of the alternating quotient kernel}

\begin{theorem}[D-TD21]
Assume the actual primitive coefficient source and complete-decomposition
hypotheses of D-TD16. Suppose its largest alternating natural degree
at any layer is \(n\ge81\), so its whole-module support parameter
is \(M=n-1\). Let \(W\) be an actual factorwise Heisenberg
coordinate of this source, with full external partner retained, and
suppose its conjugation image \(H\le\operatorname{Sp}(Q)\)
acts irreducibly on \(Q\) over \(\mathbf F_p\). Suppose there is
an actual surjection
\[
             \varphi:H\longrightarrow T\in\{A_n,S_n\},
                    \qquad N=\ker\varphi.
 \tag{TDI.9}
\]
If \(s\) is the number of \(N\)-isotypic components of \(Q\),
and \(b\) is their dual-pair block count in D-TD20, then
\[
                   b\le\lfloor(n-1)/2\rfloor,
                            \qquad s\le2b\le n-1.
 \tag{TDI.10}
\]
For \(T=A_n\), \(Q|_N\) is homogeneous. For \(T=S_n\), it
is either homogeneous or has exactly two isotypic components,
which are fixed individually by \(\varphi^{-1}(A_n)\) and exchanged
by every element with odd image. In the latter case the two
components are either a dual pair of totally isotropic spaces,
or are both nondegenerate self-dual components.

In particular this applies to the actual Heisenberg carrier in
D-TD19, with the same original complete decomposition and the same
maximal \(n\). It preserves that theorem's small-prime bound and
its faithful-original-partner scope, without strengthening either.
\end{theorem}

\begin{proof}
D-TD20 supplies an actual transitive orbit of \(b\) full complex
Heisenberg tensor factor algebras inside \(W\), all of equal local
dimension at least two. The original coefficient group maps onto
its indicated projective image, so this orbit is transitive for
that original group as well. Retain the entire complement of \(W\)
in \(V\). D-TD16 applies to this actual orbit even if it is a new
refinement rather than an orbit in the chosen recursive decomposition:
its support parameter is still computed from the fixed complete
decomposition of the whole original module. It gives the first
inequality of (TDI.10); the elementary count in (TDI.3) gives the rest.
No new maximum of alternating degrees is introduced.

The action on the \(s\) isotypic components kills \(N\), so it
is an actual homomorphism
\[
                              T\longrightarrow S_s.
 \tag{TDI.11}
\]
Since \(s\le n-1\), its restriction to \(A_n\) is trivial:
otherwise simplicity would embed \(A_n\) in \(S_s\), whereas
\(s!\le(n-1)!<n!/2=|A_n|\). If \(T=A_n\), (TDI.11)
is therefore trivial. Transitivity from D-TD20 forces \(s=1\),
which is homogeneity over \(\mathbf F_p\).

If \(T=S_n\), (TDI.11) factors through \(S_n/A_n\simeq C_2\).
A transitive action of this quotient has either one or two points.
When it has two, its nontrivial element exchanges them, so precisely
the odd images exchange the two components. Duality on a two-element
set either fixes both elements or interchanges them. The pairing
statements in D-TD20 give exactly the two alternatives asserted.
In the self-dual alternative each component has even dimension;
no further restriction on the multiplicity or on \(m\) is inferred.
\end{proof}

\paragraph{Exact remaining boundary.}
This proof is over \(\mathbf F_p\); homogeneous restriction is not
an assertion of absolute homogeneity after extension of scalars.
The modular quotient kernel \(N\) in (TDI.9) is distinct from the
kernel of the original characteristic-zero linear partner \(U\).
Nothing here treats the separate nontrivial transitive-\(\ker U\)
branch which D-TD19 excluded from its hypotheses.

The theorem does not assert semisimplicity of arbitrary modular
restrictions, a splitting of the Heisenberg normalizer, a subgroup
lift of \(A_n\) or \(S_n\), a bound on the homogeneous multiplicity,
a bound on \(m\), or a fixed-space/rank comparison with the natural
\(n\)-letter action. No actual source or target polarization is
constructed. The conclusion removes a kernel-isotypic permutation
possibility while leaving the homogeneous modular carrier as a
separate remaining problem.
